\scanpage{108}
$c\Longrightarrow a)\Longleftrightarrow b)$ trivial, pour la réciproque $a\Longrightarrow c$, noter que $E$ et $X-E=F$ étant ouverts dans $X^{\mathrm{cons}}$, sont réunions \deleted{d'ouverts} \deleted{\unsure{d'ens}} filtrantes d'ens constructibles $E_i$ resp. $F_j$, et $X$ est réunion filtrante des $E_i\cup F_j$, donc est égal à l'un des $E_i\cup F_j$, d'où $E=E_i$, \deleted{$\unsure{E\subset F_j}$} $F=F_j$ pour un $i,j$, donc $\unsure{E_i,F_j|E,F}$ constructibles.

\uline{Proposition 2}\quad Soit $f:X\longrightarrow Y$ une application continue d'espaces topologiques, telle que \deleted{\unsure{si}} pour \unsure{$\forall$} couple d'ouverts $V\subset U\subset Y$, \deleted{rétrocompact dans $Y$}\inserted{$V$ rétrocompact dans $U$}, $f^{-1}(\unsure{U|V})$ soit rétrocompact dans $f^{-1}(U)$. Alors

\margin{hyp. \unsure{trop restrictive} pour un \unsure{morphisme} de \unsure{préschémas}}

\deleted{(0) Pour \unsure{$\forall$} \unsure{immers} $Y'$ de $Y$ posons $X'=f^{-1}(Y')$, $f':X'\longrightarrow Y'$ satisfait \unsure{les mêmes hyp}.}

\unsure{(i)} Pour \unsure{$\forall$} partie constructible [loc. constructible] \deleted{\unsure{en}} $\unsure{E}$ de $Y$, $f^{-1}(\unsure{E})$ est une partie constructible [loc. constructible].

(ii) $f$ est continue de $X^{\mathrm{cons}}$ en $Y^{\mathrm{cons}}$, i.e. l'image inverse d'une partie \inserted{\deleted{loc.}} \unsure{antifique|antifinie} de $Y$ est une partie \inserted{loc.} \unsure{antifique|antifinie} de $X$.

\margin{$Y'\subset Y$}

\inserted{\unsure{car}} (i) est trivial pour ``constructible''. Pour loc. constructible, on utilise qu'on aura $Y=\bigcup Y_i$, les $Y_i$ ouverts, les $E_i=E\cap Y_i$ constructibles \inserted{dans $Y_i$}, Alors $X_i=f^{-1}(Y_i)$, \inserted{$f_i:X_i\longrightarrow Y_i$ induit,} il suffit de prouver que $f_i^{-1}(E_i)$ sont constructibles dans $X_i$, donc de prouver que $X_i\underset{\inserted{f_i}}{\unsure{-|\longrightarrow}}Y_i$ satisfait l'hyp. de $f:X\longrightarrow Y$, \deleted{c'est} Cl.

ii) Résulte trivialt de i).

\margin{107}
